\section{Exploratory results}
\label{sec:explore}

The experiments in this section were designed after the results above had been seen,
to test explanations we had offered for them; none was among our predictions
(\S\ref{sec:method}). They ran in two later batches on the same laptop and, for
\S\ref{sec:realdisk}, on a second machine.

\subsection{The interpreter, run directly}
\label{sec:gilpy}

\S\ref{sec:gil} rests on reading CPython's source: 3.12 and 3.13 hold the GIL during
an \texttt{ioctl} whose buffer exceeds 1{,}024 bytes, and 3.14 does not. We re-ran E8
unchanged under CPython 3.14.10 and under its free-threaded build 3.14t (no GIL), both
from the deadsnakes packages (Table~\ref{tab:gilpy}). Under 3.14 the fork's own lookup
and the 1{,}824-byte request gain 13--15$\times$ with 16 threads on both devices that
serve requests concurrently, nearly as much as the 984-byte request and as processes;
under 3.12 they gain nothing. Without a GIL, 16 threads gain 3.8$\times$ even on the
plain device, where every read is fast. On the one-request-at-a-time device no interpreter gains more
than 1.11$\times$. The request size therefore mattered only through the interpreter
lock, and only on 3.12 and 3.13.

\begin{table}[t]
\caption{E8 under three interpreters: speedup of 16 workers over one for FIEMAP calls
that each wait for one device read (medians of five passes).}
\label{tab:gilpy}
\small
\inputtable{../results/tab_gilpy.tex}
\end{table}

\subsection{Where each lookup waits}
\label{sec:calls}

E4 timed the fork's lookup as a whole. E9 runs each of its calls as its own pass over
500 files (\texttt{stat} all, \texttt{open} all, FIEMAP all, \texttt{close} all) and
counts device reads per pass. With 32 extents and a cold cache, the one read per file
happens inside FIEMAP on ext4 and XFS (1.00 reads per file; \texttt{stat} reads 0.01 and
0.04), and 16 threads do not make that pass faster (312 against 326\,\textmu s per file
on ext4). On F2FS the read happens in \texttt{stat} (1.02 per file, at every extent
count), and 16 threads make the \texttt{stat} pass 3.5--4.7$\times$ faster: this
confirms the attribution of \S\ref{sec:fiemap}. On Btrfs the reads split between the
two (0.18 and 0.08 per file).

F2FS's FIEMAP itself costs 350--370\,\textmu s per call, warm or cold, with no device
read, whether it asks for the whole address space or only up to the end of the file. A
BPF probe (E9b) counts 878 calls to \texttt{f2fs\_map\_blocks} per FIEMAP for files of
one and of 32 blocks alike: the walk does not depend on the file, which is why the cost
is nearly the same for one to 32 extents in Table~\ref{tab:fiemap}.

On every file system, 16 threads turn a warm \texttt{stat}, \texttt{open} or
\texttt{close} of a few microseconds into about 60\,\textmu s per file, most likely the
cost of passing the interpreter lock between 16 threads that each hold it only briefly.
The same overhead makes threads slower than one wherever there is no wait to overlap.

\subsection{Where threads could pay: the tree walk}
\label{sec:walk}

In the scattered case all device time is in the walk (\S\ref{sec:decompose}). E10
replaces the fork's single \texttt{find} with an \texttt{os.scandir} walk whose
top-level directories are split among threads; \texttt{opendir} and \texttt{readdir}
release the GIL, and the inode map came out identical to \texttt{find}'s in every run.
Under dm-delay the walk of 100{,}000 files takes 9.7\,s on one thread, 2.5\,s on four
(3.9$\times$) and 4.4\,s on sixteen; with 32 requests in flight on NBD the figures are
7.1, 1.8 and 4.6\,s. On the one-request-at-a-time device four or sixteen threads gain
1.06$\times$, and on the plain loop device they are 2.6 and 7.4 times slower than one.
Threads help the walk only when the device serves concurrent requests, and more of
them than the walk can keep busy hurt. This is the part of the analysis the fork did
not parallelise.

\subsection{Inode readahead is the kernel's own parallelism}
\label{sec:readahead}

E11 repeats the cold runs under dm-delay with \texttt{inode\_readahead\_blks} from 0
to 256 (the default is 32; the value in effect was read back from \texttt{sysfs}). In
the scattered case, readahead reads 3.6 times more data than none (7{,}287 reads
against 2{,}023 for the walk) and still makes the walk 1.7$\times$ faster (9.8 against
16.2\,s), because the kernel issues those reads together. For files created together,
the lookups read 75 blocks without readahead and 43 with the default, and 16 threads
never gain more than 1.08$\times$ at any setting, probably because threads that each
take every sixteenth file wait for the same inode-table block at the same time. The
``nothing to overlap'' of \S\ref{sec:models} does not hinge on this tunable.

\subsection{A real network disk}
\label{sec:realdisk}

Finally we ran a subset of the experiments on an OpenStack virtual machine whose disk
is network block storage (8 vCPUs, Linux 7.0, CPython 3.14.4), while its usual services
ran. Because it serves other programs we did no global cache drop: the file system under
test is unmounted before every cold run, which empties its caches. Random 4\,KiB reads
over a 2\,GiB region written just before took 35--52\,ms
at the median and up to 1.9\,s at the 99th percentile, at 6--12 reads per second one at
a time and 66--137 with 16 in flight, and the rate doubled between the start and the end
of the run. Small working sets written just before came back far faster, apparently
from a cache below the guest, which made the per-file FIEMAP passes too erratic to
interpret (16-worker intervals from 0.3$\times$ to 58$\times$). The end-to-end results agree with the laptop's:
with 1{,}000 traced files in a tree of 10{,}000 and a cold cache, upstream takes
11.6\,s, the best parallel setting 2.42\,s and the rewrite 0.29\,s; the fork's engine
gains 1.23$\times$ with 16 threads in the scattered case (8.35 against 6.79\,s); and the
threaded walk of \S\ref{sec:walk} is 6.1$\times$ faster on 16 threads than on one and
9.5$\times$ faster than \texttt{find} (1.17 against 11.1\,s), the largest gain from
threads that we measured on a real device.
